\documentclass[10pt,a4paper]{amsart}
\usepackage{amsmath,amssymb,amsthm,graphicx}
\usepackage{verbatim}
\newenvironment{reference}{\comment}{\endcomment}
%\newenvironment{reference}{}{}
\newenvironment{slogan}{\comment}{\endcomment}
\newenvironment{history}{\comment}{\endcomment}
% For commutative diagrams we use Xy-pic
\usepackage[all]{xy}

% We use 2cell for 2-commutative diagrams.
\xyoption{2cell}
% Theorem environments.
%
\theoremstyle{plain}
\newtheorem{theorem}[subsection]{Theorem}
\newtheorem{proposition}[subsection]{Proposition}
\newtheorem{lemma}[subsection]{Lemma}

\theoremstyle{definition}
\newtheorem{definition}[subsection]{Definition}
\newtheorem{example}[subsection]{Example}
\newtheorem{exercise}[subsection]{Exercise}
\newtheorem{situation}[subsection]{Situation}

\theoremstyle{remark}
\newtheorem{remark}[subsection]{Remark}
\newtheorem{remarks}[subsection]{Remarks}

\numberwithin{equation}{subsection}

% Macros
%
\def\lim{\mathop{\mathrm{lim}}\nolimits}
\def\colim{\mathop{\mathrm{colim}}\nolimits}
\def\Spec{\mathop{\mathrm{Spec}}}
\def\Hom{\mathop{\mathrm{Hom}}\nolimits}
\def\Ext{\mathop{\mathrm{Ext}}\nolimits}
\def\SheafHom{\mathop{\mathcal{H}\!\mathit{om}}\nolimits}
\def\SheafExt{\mathop{\mathcal{E}\!\mathit{xt}}\nolimits}
\def\Sch{\mathit{Sch}}
\def\Mor{\mathop{\mathrm{Mor}}\nolimits}
\def\Ob{\mathop{\mathrm{Ob}}\nolimits}
\def\Sh{\mathop{\mathit{Sh}}\nolimits}
\def\NL{\mathop{N\!L}\nolimits}
\def\CH{\mathop{\mathrm{CH}}\nolimits}
\def\proetale{{pro\text{-}\acute{e}tale}}
\def\etale{{\acute{e}tale}}
\def\QCoh{\mathit{QCoh}}
\def\Ker{\mathop{\mathrm{Ker}}}
\def\Im{\mathop{\mathrm{Im}}}
\def\Coker{\mathop{\mathrm{Coker}}}
\def\Coim{\mathop{\mathrm{Coim}}}

% Boxtimes
%
\DeclareMathSymbol{\boxtimes}{\mathbin}{AMSa}{"02}

%
% Macros for moduli stacks/spaces
%
\def\QCohstack{\mathcal{QC}\!\mathit{oh}}
\def\Cohstack{\mathcal{C}\!\mathit{oh}}
\def\Spacesstack{\mathcal{S}\!\mathit{paces}}
\def\Quotfunctor{\mathrm{Quot}}
\def\Hilbfunctor{\mathrm{Hilb}}
\def\Curvesstack{\mathcal{C}\!\mathit{urves}}
\def\Polarizedstack{\mathcal{P}\!\mathit{olarized}}
\def\Complexesstack{\mathcal{C}\!\mathit{omplexes}}
% \Pic is the operator that assigns to X its picard group, usage \Pic(X)
% \Picardstack_{X/B} denotes the Picard stack of X over B
% \Picardfunctor_{X/B} denotes the Picard functor of X over B
\def\Pic{\mathop{\mathrm{Pic}}\nolimits}
\def\Picardstack{\mathcal{P}\!\mathit{ic}}
\def\Picardfunctor{\mathrm{Pic}}
\def\Deformationcategory{\mathcal{D}\!\mathit{ef}}
\usepackage{cleveref}
\crefname{originalequation}{original equation}{original equations}
\Crefname{originalequation}{Original equation}{Original equations}
\makeatletter
\newlabel{equation-tame-symbol}{{0EAQ}{source}{}{sourcepointer-1}{}}
\newlabel{equation-tame-symbol@cref}{{[originalequation][0][]0EAQ}{[1][0][]source}{}{}{}}
\makeatother
\crefname{originalreference}{original reference}{original references}
\Crefname{originalreference}{Original reference}{Original references}
\makeatletter
\newlabel{item-1-x}{{0EAK}{source}{}{sourcepointer-2}{}}
\newlabel{item-1-x@cref}{{[originalreference][0][]0EAK}{[1][0][]source}{}{}{}}
\makeatother
\crefname{originalreference}{original reference}{original references}
\Crefname{originalreference}{Original reference}{Original references}
\makeatletter
\newlabel{item-1-x-better}{{0EAP}{source}{}{sourcepointer-3}{}}
\newlabel{item-1-x-better@cref}{{[originalreference][0][]0EAP}{[1][0][]source}{}{}{}}
\makeatother
\crefname{originalreference}{original reference}{original references}
\Crefname{originalreference}{Original reference}{Original references}
\makeatletter
\newlabel{item-bilinear}{{0EAI}{source}{}{sourcepointer-4}{}}
\newlabel{item-bilinear@cref}{{[originalreference][0][]0EAI}{[1][0][]source}{}{}{}}
\makeatother
\crefname{originalreference}{original reference}{original references}
\Crefname{originalreference}{Original reference}{Original references}
\makeatletter
\newlabel{item-bilinear-better}{{0EAL}{source}{}{sourcepointer-5}{}}
\newlabel{item-bilinear-better@cref}{{[originalreference][0][]0EAL}{[1][0][]source}{}{}{}}
\makeatother
\crefname{originalreference}{original reference}{original references}
\Crefname{originalreference}{Original reference}{Original references}
\makeatletter
\newlabel{item-normalization}{{0EAN}{source}{}{sourcepointer-6}{}}
\newlabel{item-normalization@cref}{{[originalreference][0][]0EAN}{[1][0][]source}{}{}{}}
\makeatother
\crefname{originalreference}{original reference}{original references}
\Crefname{originalreference}{Original reference}{Original references}
\makeatletter
\newlabel{item-skew}{{0EAJ}{source}{}{sourcepointer-7}{}}
\newlabel{item-skew@cref}{{[originalreference][0][]0EAJ}{[1][0][]source}{}{}{}}
\makeatother
\crefname{originalreference}{original reference}{original references}
\Crefname{originalreference}{Original reference}{Original references}
\makeatletter
\newlabel{item-skew-better}{{0EAM}{source}{}{sourcepointer-8}{}}
\newlabel{item-skew-better@cref}{{[originalreference][0][]0EAM}{[1][0][]source}{}{}{}}
\makeatother
\crefname{originallemma}{original lemma}{original lemmas}
\Crefname{originallemma}{Original lemma}{Original lemmas}
\makeatletter
\newlabel{lemma-prepare-tame-symbol}{{0EAG}{source}{}{sourcepointer-9}{}}
\newlabel{lemma-prepare-tame-symbol@cref}{{[originallemma][0][]0EAG}{[1][0][]source}{}{}{}}
\makeatother
\title{Bilinearity, skew symmetry and the two normalization properties of the tame symbol}
\author{The Stacks Project Authors}
\date{stacks-project:chow.tex:local-snapshot-20261012; GNU FDL 1.2 or later}
\begin{document}
\maketitle
\noindent\textit{Editorial note.} All delivered mathematical text is the author's own, copied from the pinned Stacks Project chapter file chow.tex; no statement or proof step was rewritten, reordered or removed.\par\medskip
\noindent\textit{Editorial note.} Preamble provenance: the pinned chapter file opens with \textbackslash{}input\{preamble\} and carries no preamble text of its own, so the project's single Stacks-wide preamble staged at sources/stacks-chow/preamble.tex is excerpted verbatim as three spans, lines 8-12 (the verbatim package together with the project's own reference, slogan and history display-environment definitions, including the author's commented-out older form of the reference environment), lines 14-18 (Xy-pic setup) and lines 154-217 (the native theorem declarations and mathematical macros). The lines 8-12 span is required because a counted statement or its surrounding text may use those display environments, which the project typesets as comments. Nothing was added, removed or redefined; the wrapper class and the amsmath/amssymb/amsthm/graphicx base packages are supplied by the assembler.\par\medskip
\noindent\textit{Editorial note.} Same-chapter pointers are located at their \textbackslash{}label lines inside the pinned chapter file; cross-chapter pointers are located in the project's own tag list staged in the snapshot (sources/stacks-chow/tags.txt), at the verbatim lines carrying their canonical tags.\par\medskip
\noindent\textit{Allowed standard prerequisites (editorial).}\begin{itemize}
\item The preliminary preparation of the tame symbol in a finite extension. Same chapter chow.tex, label lemma-prepare-tame-symbol, invoked at proof lines 913 and 933. Produces, for the finitely many nonzerodivisors involved, a finite ring extension in which the elements become monomials in a single nonzerodivisor up to units, which reduces each verification to an identity among the units and the exponents.
\end{itemize}\medskip
\ifdefined\hypertarget\hypertarget{sourcepointer-1}{}\fi
\noindent\textit{Source reference.} Original equation 0EAQ: sources/stacks-chow/pinned-original-source.tex, source lines 820--820. Navigation pointer for \textbackslash{}ref\{equation-tame-symbol\} used inside the counted span (line 903); the argument that consumes it is retained in the span.\par\medskip
\ifdefined\hypertarget\hypertarget{sourcepointer-2}{}\fi
\noindent\textit{Source reference.} Original reference 0EAK: sources/stacks-chow/pinned-original-source.tex, source lines 779--779. Navigation pointer for \textbackslash{}ref\{item-1-x\} used inside the counted span (line 907); the argument that consumes it is retained in the span.\par\medskip
\ifdefined\hypertarget\hypertarget{sourcepointer-3}{}\fi
\noindent\textit{Source reference.} Original reference 0EAP: sources/stacks-chow/pinned-original-source.tex, source lines 794--794. Navigation pointer for \textbackslash{}ref\{item-1-x-better\} used inside the counted span (line 905); the argument that consumes it is retained in the span.\par\medskip
\ifdefined\hypertarget\hypertarget{sourcepointer-4}{}\fi
\noindent\textit{Source reference.} Original reference 0EAI: sources/stacks-chow/pinned-original-source.tex, source lines 773--773. Navigation pointer for \textbackslash{}ref\{item-bilinear\} used inside the counted span (line 907); the argument that consumes it is retained in the span.\par\medskip
\ifdefined\hypertarget\hypertarget{sourcepointer-5}{}\fi
\noindent\textit{Source reference.} Original reference 0EAL: sources/stacks-chow/pinned-original-source.tex, source lines 782--782. Navigation pointer for \textbackslash{}ref\{item-bilinear-better\} used inside the counted span (line 904); the argument that consumes it is retained in the span.\par\medskip
\ifdefined\hypertarget\hypertarget{sourcepointer-6}{}\fi
\noindent\textit{Source reference.} Original reference 0EAN: sources/stacks-chow/pinned-original-source.tex, source lines 790--790. Navigation pointer for \textbackslash{}ref\{item-normalization\} used inside the counted span (line 905); the argument that consumes it is retained in the span.\par\medskip
\ifdefined\hypertarget\hypertarget{sourcepointer-7}{}\fi
\noindent\textit{Source reference.} Original reference 0EAJ: sources/stacks-chow/pinned-original-source.tex, source lines 776--776. Navigation pointer for \textbackslash{}ref\{item-skew\} used inside the counted span (line 907); the argument that consumes it is retained in the span.\par\medskip
\ifdefined\hypertarget\hypertarget{sourcepointer-8}{}\fi
\noindent\textit{Source reference.} Original reference 0EAM: sources/stacks-chow/pinned-original-source.tex, source lines 786--786. Navigation pointer for \textbackslash{}ref\{item-skew-better\} used inside the counted span (line 904); the argument that consumes it is retained in the span.\par\medskip
\ifdefined\hypertarget\hypertarget{sourcepointer-9}{}\fi
\noindent\textit{Source reference.} Original lemma 0EAG: sources/stacks-chow/pinned-original-source.tex, source lines 642--642. Navigation pointer for \textbackslash{}ref\{lemma-prepare-tame-symbol\} used inside the counted span (line 913); the argument that consumes it is retained in the span.\par\medskip
\begin{lemma}
\label{lemma-tame-symbol}
The tame symbol (\ref{equation-tame-symbol}) satisfies
(\ref{item-bilinear-better}), (\ref{item-skew-better}),
(\ref{item-normalization}), (\ref{item-1-x-better}) and hence
gives a map $\partial_A : Q(A)^* \times Q(A)^* \to \kappa(\mathfrak m)^*$
satisfying (\ref{item-bilinear}), (\ref{item-skew}), (\ref{item-1-x}).
\end{lemma}

\begin{proof}
Let us prove (\ref{item-bilinear-better}).
Let $a_1, a_2, a_3 \in A$ be nonzerodivisors.
Choose $A \subset B$ as in Lemma \ref{lemma-prepare-tame-symbol}
for $a_1, a_2, a_3$. Then the equality
$$
\partial_A(a_1a_2, a_3) = \partial_A(a_1, a_3) \partial_A(a_2, a_3)
$$
follows from the equality
$$
(-1)^{(e_{1, j} + e_{2, j})e_{3, j}}
(u_{1, j}u_{2, j})^{e_{3, j}}u_{3, j}^{-e_{1, j} - e_{2, j}} =
(-1)^{e_{1, j}e_{3, j}}
u_{1, j}^{e_{3, j}}u_{3, j}^{-e_{1, j}}
(-1)^{e_{2, j}e_{3, j}}
u_{2, j}^{e_{3, j}}u_{3, j}^{-e_{2, j}}
$$
in $B_j$. Properties (\ref{item-skew-better}) and
(\ref{item-normalization}) are equally immediate.

\medskip\noindent
Let us prove (\ref{item-1-x-better}). Let $a_1, a_2, a_1 - a_2 \in A$
be nonzerodivisors and set $a_3 = a_1 - a_2$.
Choose $A \subset B$ as in Lemma \ref{lemma-prepare-tame-symbol}
for $a_1, a_2, a_3$. Then it suffices to show
$$
(-1)^{e_{1, j}e_{2, j} + e_{1, j}e_{3, j} + e_{2, j}e_{3, j} + e_{2, j}}
u_{1, j}^{e_{2, j} - e_{3, j}}
u_{2, j}^{e_{3, j} - e_{1, j}}
u_{3, j}^{e_{1, j} - e_{2, j}} \bmod \mathfrak m_j = 1
$$
This is clear if $e_{1, j} = e_{2, j} = e_{3, j}$.
Say $e_{1, j} > e_{2, j}$. Then we see that $e_{3, j} = e_{2, j}$
because $a_3 = a_1 - a_2$ and we see that $u_{3, j}$
has the same residue class as $-u_{2, j}$. Hence
the formula is true -- the signs work out as well
and this verification is the reason for the choice of signs
in (\ref{equation-tame-symbol}).
The other cases are handled in exactly the same manner.
\end{proof}

\end{document}
